\section{STAIR: Manipulating Collaborative and Multimodal Information}

\subsection{Kết quả tái lập thực nghiệm gốc}
\label{sec:stair_reproduction}

\subsubsection{Lý do thực hiện tái lập}

Trước khi tiếp tục các phương án cải tiến, nhóm nhận thấy lần thực nghiệm baseline đầu tiên còn một số điểm cấu hình chưa đúng theo mô tả của tác giả, dẫn đến việc lựa chọn checkpoint tốt nhất chưa nhất quán với quy trình đánh giá của paper. Để đảm bảo mọi so sánh ở các phần cải tiến tiếp theo được thực hiện trên một nền tảng baseline đáng tin cậy, nhóm tiến hành chạy lại toàn bộ thực nghiệm gốc của STAIR với cấu hình đã sửa đúng, cụ thể là chỉ số dùng để chọn checkpoint tốt nhất trên tập validation là NDCG@20, cùng bộ theo dõi đầy đủ gồm Recall@10, Recall@20, NDCG@10 và NDCG@20. Đây là bước tái lập độc lập, đóng vai trò xác nhận độ tin cậy của pipeline trước khi dùng làm điểm xuất phát cho các cải tiến ở đợt 2.

\subsubsection{Thiết lập thực nghiệm}

Quá trình tái lập được thực hiện trên ba tập dữ liệu Amazon2014, gồm Baby, Sports và Electronics, mỗi tập huấn luyện tối đa 500 epoch. Checkpoint tốt nhất được chọn dựa trên giá trị NDCG@20 cao nhất đo được trên tập validation trong suốt quá trình huấn luyện, sau đó mô hình tại checkpoint này được đánh giá lại trên tập test để lấy kết quả cuối cùng. Toàn bộ số liệu công bố trong phần này đều là kết quả trên tập test tại đúng checkpoint đã chọn, giữ nguyên quy ước đánh giá đã áp dụng xuyên suốt các phần thực nghiệm trước đó của khóa luận.

\subsubsection{Kết quả so sánh với paper gốc}

Bảng~\ref{tab:stair_reproduction} tổng hợp kết quả tái lập trên cả ba tập dữ liệu, đối chiếu trực tiếp với số liệu công bố trong Table 2 của paper STAIR.

\renewcommand{\arraystretch}{1.15}
\begingroup
\footnotesize
\setlength{\tabcolsep}{5pt}
\begin{longtable}{%
    >{\raggedright\arraybackslash}p{1.9cm}
    >{\centering\arraybackslash}p{1.6cm}
    >{\centering\arraybackslash}p{1.9cm}
    >{\centering\arraybackslash}p{1.9cm}
    >{\centering\arraybackslash}p{1.9cm}}

    \caption{Kết quả tái lập thực nghiệm STAIR so với paper gốc trên ba tập dữ liệu.}
    \label{tab:stair_reproduction} \\
    \hline
    \rowcolor{gray!15}\textbf{Tập dữ liệu} & \textbf{Chỉ số} & \textbf{Kết quả tái lập} & \textbf{Paper (Table 2)} & \textbf{Chênh lệch} \\
    \hline
    \endfirsthead
    \multicolumn{5}{c}{{\bfseries Bảng \thetable\ (tiếp theo)}} \\
    \hline
    \rowcolor{gray!15}\textbf{Tập dữ liệu} & \textbf{Chỉ số} & \textbf{Kết quả tái lập} & \textbf{Paper (Table 2)} & \textbf{Chênh lệch} \\
    \hline
    \endhead
    \midrule
    \multicolumn{5}{r}{\textit{Tiếp tục ở trang sau...}} \\
    \endfoot
    \hline
    \endlastfoot

    Baby & Recall@10 & 0.0674 & 0.0674 & $0.00\%$ \\
    \midrule
    {(best epoch 455/500)} & Recall@20 & 0.1042 & 0.1042 & $0.00\%$ \\
    \midrule
    & NDCG@10   & 0.0359 & 0.0359 & $0.00\%$ \\
    \midrule
    & NDCG@20   & 0.0454 & 0.0453 & $+0.22\%$ \\
    \hline
    Sports & Recall@10 & 0.0743 & 0.0743 & $0.00\%$ \\
    \midrule
    {(best epoch 500/500)} & Recall@20 & 0.1111 & 0.1117 & $-0.54\%$ \\
    \midrule
    & NDCG@10   & 0.0405 & 0.0407 & $-0.49\%$ \\
    \midrule
    & NDCG@20   & 0.0500 & 0.0503 & $-0.60\%$ \\
    \hline
    Electronics & Recall@10 & 0.0442 & 0.0440 & $+0.45\%$ \\
    \midrule
    {(best epoch 490/500)} & Recall@20 & 0.0665 & 0.0663 & $+0.30\%$ \\
    \midrule
    & NDCG@10   & 0.0246 & 0.0245 & $+0.41\%$ \\
    \midrule
    & NDCG@20   & 0.0303 & 0.0302 & $+0.33\%$ \\
    \hline
\end{longtable}
\endgroup

Trên toàn bộ mười hai chỉ số đo được, tức bốn chỉ số trên mỗi tập trong số ba tập dữ liệu, chênh lệch tuyệt đối so với paper đều nằm dưới 1\%. Đây là mức sai lệch nằm trong biên độ nhiễu tự nhiên giữa một lần chạy đơn lẻ so với kết quả trung bình của nhiều lần chạy mà tác giả công bố, không phản ánh lỗi trong pipeline tái lập.

\subsubsection{Phân tích chi tiết theo từng tập dữ liệu}

\paragraph{Tập Baby — khớp gần như tuyệt đối:}
Kết quả trên Baby trùng khớp hoàn toàn với paper ở ba trong bốn chỉ số, chỉ lệch nhẹ $+0.22\%$ ở NDCG@20. Checkpoint tốt nhất rơi vào epoch 455 trên tổng số 500 epoch huấn luyện, tức mô hình đã đạt trạng thái ổn định trước khi kết thúc quá trình huấn luyện khoảng 45 epoch. Điều này cho thấy 500 epoch là đủ để mô hình hội tụ hoàn toàn trên tập dữ liệu này, và kết quả thu được phản ánh đúng khả năng tối đa của mô hình trong cấu hình đã cho.

\paragraph{Tập Sports — lệch nhẹ dưới 1\%, gắn liền với việc chưa hội tụ hoàn toàn:}
Sports là tập duy nhất có kết quả thấp hơn paper một chút, dao động từ $-0.49\%$ đến $-0.60\%$ trên ba chỉ số Recall@20, NDCG@10 và NDCG@20. Đây cũng là tập duy nhất có checkpoint tốt nhất rơi đúng vào epoch cuối cùng của quá trình huấn luyện, tức epoch 500 trên 500. Việc checkpoint tốt nhất nằm ở đúng ranh giới trên cho thấy giá trị NDCG@20 trên tập validation nhiều khả năng vẫn đang trên đà tăng tại thời điểm huấn luyện dừng lại, chưa chắc đã chạm đến điểm hội tụ thật sự. Đây là cách giải thích hợp lý cho mức chênh lệch nhỏ quan sát được trên tập này, và khác biệt so với Baby, nơi checkpoint tốt nhất xuất hiện sớm hơn nhiều so với epoch cuối.

\paragraph{Tập Electronics — kết quả nhỉnh hơn paper trên toàn bộ chỉ số:}
Trên tập Electronics, kết quả tái lập cao hơn paper trên cả bốn chỉ số, với mức chênh lệch từ $+0.30\%$ đến $+0.45\%$. Đây là tập dữ liệu có quy mô lớn nhất trong ba tập, với khoảng 192 nghìn người dùng và 1.69 triệu tương tác, nên việc pipeline tái lập cho kết quả ổn định và nhỉnh hơn nhẹ trên một tập dữ liệu lớn như vậy càng củng cố thêm độ tin cậy của quy trình huấn luyện và đánh giá đã thiết lập. Checkpoint tốt nhất được chọn tại epoch 490 trên 500, cho thấy mô hình cũng đã gần đạt trạng thái hội tụ ổn định trước khi kết thúc huấn luyện.

\subsubsection{Kết luận}

Kết quả tái lập trên cả ba tập dữ liệu đều đạt độ khớp cao với số liệu công bố trong paper gốc, với sai lệch tuyệt đối dưới 1\% trên toàn bộ mười hai chỉ số đánh giá. Mức chênh lệch nhỏ quan sát được trên tập Sports có nguyên nhân hợp lý là quá trình huấn luyện chưa hoàn toàn hội tụ trong giới hạn 500 epoch, không xuất phát từ sai sót trong cấu hình hay quy trình đánh giá. Kết quả này xác nhận pipeline tái lập đã hoạt động đúng theo mô tả của tác giả và đủ tin cậy để dùng làm nền tảng so sánh cho các phương án cải tiến được trình bày ở các phần tiếp theo của đợt 2.

\clearpage